% Solar Filament Segmentation Challenge 2026 — Technical Report (4 pages)
% Paste this into the competition's Overleaf template; content fits ~4 pages.
\documentclass[10pt]{article}
\usepackage[margin=2.2cm]{geometry}
\usepackage{graphicx, booktabs, amsmath, hyperref}
\title{Attention-Gated U-Net with Skeleton-Aware Loss for\\
Solar Filament Segmentation (Panoptic Quality Optimisation)}
\author{Team Name}
\date{}

\begin{document}
\maketitle

\begin{abstract}
We address the Solar Filament Segmentation Challenge 2026 (MAGFiLO v1.0, GONG
H-$\alpha$ observations, 2048$\times$2048 px) under the Panoptic Quality (PQ)
metric. Our solution combines: (i) an anti-leakage GroupKFold protocol keyed on
physical observation names, since the same image is independently annotated by
2--3 annotators; (ii) a timm-encoder U-Net with additive attention gates and a
zero-initialised dilated bottleneck; (iii) a BCE + soft-Dice + soft-clDice loss
that penalises fragmented thin structures while using segmentation masks only
for training; (iv) sliding-window inference with 8-way dihedral TTA and fold
ensembling; (v) distance-transform watershed instance separation with
GT-style post-processing (hole filling, single connected piece, small-object
suppression) and post-processing parameters tuned directly on local PQ.
This report is a draft; quantitative comparisons remain to be completed.
\end{abstract}

\section{Task and Data Analysis}
Filaments are thin dark structures (mean $\approx$2\,k px area inside
$\approx$16\,k px bounding boxes) on full-disk H-$\alpha$ images. Three dataset
properties drive our design: (1) duplicated observations (same physical image,
2--3 independent annotator sets; COCO ids differ only by a batch prefix);
(2) the $\pm70^{\circ}$ central-meridian annotation limit; (3) inter-annotator
agreement $\kappa=0.66$, i.e.\ instance definition (one filament vs.\ several
pieces) is itself noisy, so post-merging/splitting behaviour must be controlled.

\section{Method}
\subsection{Model}
U-Net over a timm \texttt{features\_only} encoder (ResNet-18 default), additive
attention gates on skip connections~\cite{oktay2018}, and a residual dilated
bottleneck ($d\!=\!1,2,4$; GroupNorm scale zero-initialised so training starts
from the plain-U-Net function). The segmentation head is trained from polygon masks;
other ground-truth metadata is not used as a training target.

\subsection{Loss}
$\mathcal{L} = 0.3\,\mathrm{BCE} + 0.3\,\mathrm{softDice} + 0.4\,(1-\mathrm{clDice})$,
all terms masked to the annotated limb region. Soft-clDice~\cite{shit2021} compares differentiable skeletons and
penalises filament fragmentation invisible to Dice.

\subsection{Training}
Filament-centred 768$\times$768 crops (75\%), dihedral augmentations
(flips/rot90 are physically valid), photometric jitter, AMP, AdamW +
cosine schedule, weight EMA.

\subsection{Inference and Post-processing}
Overlap-add sliding window (1024 px, 256 overlap) with 8-way dihedral TTA;
probability maps averaged over folds. Instances via thresholding and connected
components by default, with optional EDT watershed, then cleanup: hole filling,
largest connected component, $\geq$50 px. Threshold, watershed distance and
min-area are grid-tuned on the held-out fold using a local PQ implementation.

\section{Experiments}
GroupKFold ($k{=}5$) keyed on observation names; validation uses one annotation
set per image and is scored with local PQ (greedy IoU matching @ 0.5).
Ablation (fold 0, local PQ): \emph{fill in table}.
\begin{center}\small
\begin{tabular}{@{}lc@{}}
\toprule
variant & local PQ \\
\midrule
U-Net, BCE+Dice & -- \\
+ clDice mask loss & -- \\
+ attention gates + dilated bottleneck & -- \\
+ TTA + fold ensemble + PQ-tuned post-proc & -- \\
\bottomrule
\end{tabular}
\end{center}

\section{Conclusion}
We hypothesise that skeleton-aware losses and instance-aware post-processing
improve PQ on thin solar structures. Controlled ablations are still required
to measure their contributions and compare them with architecture changes.

\begin{thebibliography}{9}
\bibitem{oktay2018} Oktay et al., ``Attention U-Net'', MICCAI 2018.
\bibitem{shit2021} Shit et al., ``clDice -- a Novel Connectivity-Preserving Loss
Function for Deep Learning on Tubular Structures'', CVPR 2021.
\bibitem{ahmadzadeh2024} Ahmadzadeh et al., ``A dataset of manually annotated
filaments from H-alpha observations'', \emph{Sci.\ Data}, 2024.
\bibitem{kirillov2019} Kirillov et al., ``Panoptic Segmentation'', CVPR 2019.
\end{thebibliography}

\end{document}
